\subsection{Определенные интегралы в физике}
%%%%%%%%%%%%%%%%%%%%%%%%%%%%%%%%%%%%%%%%%%%%%%%%%%%%
%%%%%%%%%%%%%%%%%%%%%%%%%%%%%%%%%%%%%%%%%%%%%%%%%%%%
 \z{Тело движется вдоль оси $x$ с постоянным ускорением $a$. Дважды интегрируя равенство $x''(t) = a$, получите зависимость скорости $v(t)$ и координаты $x(t)$ тела от времени. В начальный момент $t = 0$ скорость и координата тела равны соответственно $v_0$ и $x_0$.}
%
{$v(t) = v_0 + at;$~$x(t) = x_0 + v_0 t + \frac{at^2}{2}$}
%%%%%%%%%%%%%%%%%%%%%%%%%%%%%%%%%%%%%%%%%%%%%%%%%%%%
%%%%%%%%%%%%%%%%%%%%%%%%%%%%%%%%%%%%%%%%%%%%%%%%%%%% 
 \z{Тело движется вдоль оси $x$ с переменным ускорением $a(t) = a_0 + bt$. Дважды интегрируя равенство $x''(t) = a_0 + bt$, получите зависимость скорости $v(t)$ и координаты $x(t)$ тела от времени. В начальный момент $t = 0$ скорость и координата тела равны соответственно $v_0$ и $x_0$.}
%
{$v(t) = v_0 + a_0 t + \frac{bt^2}{2}$~$x(t) = x_0 + v_0 t + \frac{a_0 t^2}{2} + \frac{bt^3}{6}$}
%%%%%%%%%%%%%%%%%%%%%%%%%%%%%%%%%%%%%%%%%%%%%%%%%%%% 
 \z{Какое давление оказывал бы на поверхность Земли однородный железный столб, высота которого равна радиусу Земли $R$? Плотность железа $\rho = 7800\,\text{кг/м}^3$.}
%
{$P = \frac{1}{2}\rho g R \approx 2.5\cdot10^{11}\,\text{Па}$}
%%%%%%%%%%%%%%%%%%%%%%%%%%%%%%%%%%%%%%%%%%%%%%%%%%%%
\z{Лодка длины $L$, двигаясь по инерции, наезжает на берег и останавливается из-за трения, когда половина её длины оказывается на суше. Найдите начальную скорость лодки $v_0$. Коэффициент трения равен $\mu$.}
%
{$v_0 =\sqrt{\frac{\mu gL}{4}}$}
%%%%%%%%%%%%%%%%%%%%%%%%%%%%%%%%%%%%%%%%%%%%%%%%%%%%	
\z{Где находится центр масс полушара радиуса $R$?}
%
{на расстоянии $\frac{3}{8}R$ от центра}
